\par\penalty0
\section{Tái định nghĩa ngăn xếp hệ thống (figure 2 mới) và khái niệm ranh giới thuê bao}
\label{outline:1:reference-model:1}

\par\penalty0
\subsection{Phân tích Ngăn xếp Phần mềm Nguyên bản (Figure 2) của Kabbedijk et al. (2015)}
\label{outline:1:reference-model:2}

\par\penalty0
\subsubsection{Ý nghĩa của quy ước nét liền và nét đứt trong Figure 2 nguyên bản}
\label{outline:1:reference-model:3}

\par\penalty0
\subsubsection{Giá trị và giới hạn của mô hình đối với nghiên cứu hiện tại}
\label{outline:1:reference-model:4}

\par\penalty0
\subsection{Động lực Dịch chuyển Công nghệ và Đề xuất Mô hình Ngăn xếp Cập nhật cho Kỷ nguyên Cloud-Native}
\label{outline:1:reference-model:5}

\par\penalty0
\subsubsection{Động lực 1: Sự thay đổi vai trò của Middleware trong kiến trúc đa thuê bao hiện đại}
\label{outline:1:reference-model:6}

\par\penalty0
\subsubsection{Động lực 2: Sự phổ biến của Container và cơ chế cô lập ở cấp nhân hệ điều hành}
\label{outline:1:reference-model:7}

\par\penalty0
\subsubsection{Quy ước ký hiệu trong mô hình ngăn xếp cập nhật}
\label{outline:1:reference-model:8}

\par\penalty0
\subsection{Định nghĩa Chi tiết 9 Tầng Thành phần trong Ngăn xếp Đa thuê bao Hiện đại}
\label{outline:1:reference-model:9}

\par\penalty0
\subsubsection{Hardware (Phần cứng)}
\label{outline:1:reference-model:10}

\par\penalty0
\subsubsection{Virtual Machine (Máy ảo)}
\label{outline:1:reference-model:11}

\par\penalty0
\subsubsection{Operating System (Hệ điều hành)}
\label{outline:1:reference-model:12}

\par\penalty0
\subsubsection{Database Server (Máy chủ Cơ sở dữ liệu)}
\label{outline:1:reference-model:13}

\par\penalty0
\subsubsection{Database (Cơ sở dữ liệu)}
\label{outline:1:reference-model:14}

\par\penalty0
\subsubsection{Database Schema (Lược đồ Cơ sở dữ liệu)}
\label{outline:1:reference-model:15}

\par\penalty0
\subsubsection{Middleware (Phần mềm trung gian)}
\label{outline:1:reference-model:16}

\par\penalty0
\subsubsection{Application Server (Máy chủ Ứng dụng)}
\label{outline:1:reference-model:17}

\par\penalty0
\subsubsection{Application Instance (Thực thể Ứng dụng)}
\label{outline:1:reference-model:18}

\par\penalty0
\subsection{Đề xuất Khái niệm “Ranh giới Thuê bao” (Tenant Boundary) từ Nghịch lý của Tính Trong suốt}
\label{outline:1:reference-model:19}

\par\penalty0
\subsubsection{Nghịch lý của tính trong suốt trong kiến trúc đa thuê bao}
\label{outline:1:reference-model:20}

\par\penalty0
\subsubsection{Đề xuất định nghĩa Ranh giới Thuê bao}
\label{outline:1:reference-model:21}

\par\penalty0
\subsubsection{Ba thuộc tính bất biến của Ranh giới Thuê bao}
\label{outline:1:reference-model:22}

\par\penalty0
\subsubsection{Bốn chiều Ranh giới Thực thi}
\label{outline:1:reference-model:23}

\par\penalty0
\subsubsection{Phổ Độ dày Ranh giới}
\label{outline:1:reference-model:24}

\par\penalty0
\subsubsection{Quy luật đánh đổi kiến trúc và sự dịch chuyển trách nhiệm cô lập}
\label{outline:1:reference-model:25}

\par\penalty0
\subsection{Mô hình Hóa Tổng thể Kiến trúc Đa thuê bao (Mô hình Tích hợp Đầy đủ)}
\label{outline:1:reference-model:26}
